\begin{abstract}
有限库存使平台面临两个决策：售罄后用什么商品补足推荐位，以及是否把尚有库存的商品留给后续顾客。本文构造单类型算例，说明固定 Top-$K$ 推荐在 20 期内相对简单可行策略损失 47.08\% 的期望收入。对于课程实例，含库存约束的流体线性规划给出收入上界 706，商品 1 的影子价格为 9。本文据此设计滚动策略 \LPDual{}，每期按购买概率与扣除库存机会成本后的价格为有货商品排序。五份给定到达序列上，该策略的平均期望收入为 \ERDual，平均和最大期望遗憾分别为 \GavgDual{} 和 \GmaxDual，缺货推荐次数均为零。相对库存感知贪心，每条序列平均增收 \DualGreedyAvg{}（\DualGreedyRelativePct\%），与预知到达顺序的最优策略之间的平均差距由 \LavgGreedy{} 降至 \LavgDual{}，缩小 \DualSeqGapReductionPct\%。库存缩放实验表明，缺货展示在库存极少时最频繁，而把稀缺商品分配给不同顾客带来的收入差异在库存中度紧张时最大。
\end{abstract}
\noindent\textbf{关键词：}在线推荐；库存约束；流体线性规划；影子价格；动态规划
